\chapter{Typed Hypergraphs and Content Atoms}
\label{ch:typed-hypergraphs}

\section{Overview}

The preceding chapters developed the foundational geometric and variational
structure of the RSVP plenum and the cognitive rewrite engine.
We now introduce the second core substrate of the theory:
\emph{typed Polyxan hypergraphs}. These objects represent the explicit,
structural, content-level configurations that inhabit the semantic manifold
$\mathcal{M}$ and evolve under rewrite dynamics.

Typed hypergraphs provide:
\begin{itemize}
  \item a \emph{combinatorial skeleton} for semantic content,
  \item a \emph{fibered sheaf representation} enabling geometric embedding,
  \item a \emph{type-theoretic layer} organizing content families,
  \item and a \emph{hyperedge algebra} supporting quines, recursion, and
        structural equivalence.
\end{itemize}

The purpose of this chapter is to construct the typed hypergraph system
underlying Polyxan content atoms and prepare for the RSVP--Polyxan coupling
(Chapter~\ref{ch:coupling}).

\section{Typed Hypergraphs}

\subsection{Basic definition}

A Polyxan hypergraph is not merely a combinatorial hypergraph; it is a geometric
object fibered over $\mathcal{M}$.

\begin{definition}[Typed Polyxan Hypergraph]
A typed Polyxan hypergraph is a triple
\[
H = (V,E,\tau)
\]
where:
\begin{enumerate}
  \item $V$ is a set of vertices equipped with an embedding
        \[
        \sigma_V : V \hookrightarrow \mathcal{M},
        \]
  \item $E$ is a multiset of hyperedges, each of the form
        \[
        e = (v_1,\dots,v_k)\in V^k,
        \]
  \item $\tau$ is a type assignment
        \[
        \tau : V\amalg E \to \mathcal{T}
        \]
        where $\mathcal{T}$ is a type universe.
\end{enumerate}
\end{definition}

Embedding vertices into $\mathcal{M}$ ensures that every structural element has
semantic localization.

\subsection{Hyperedge signatures}

Each edge type $t\in\mathcal{T}$ has a signature
\[
\mathrm{sig}(t) : \mathcal{T}^k \longrightarrow \mathcal{T}
\]
determining how vertex types combine to produce edge types. This enables rich,
multi-typed combinatorial structure.

\section{Fibered Hypergraphs and Sheaf Structure}

\subsection{The hypergraph fibration}

Let
\[
\pi_H : H \longrightarrow \mathcal{M}
\]
map each vertex $v\mapsto \sigma_V(v)$ and extend to edges via their incident
vertices.

\begin{definition}[Semantic Hypergraph Fibration]
A typed hypergraph $H$ is \emph{semantic} if $\pi_H$ is a Grothendieck
fibration, i.e. for every morphism in $\mathcal{M}$ and every object in a
fiber, there exists a Cartesian lift.
\end{definition}

Intuitively: hypergraph structure must lift naturally along semantic flows.

\subsection{Hypergraph sheaf}

Define the sheaf of hypergraphs:

\begin{definition}[Polyxan Hypergraph Sheaf]
For each open $U\subseteq\mathcal{M}$, let
\[
\mathscr{H}(U)
\]
be the category of typed hypergraphs whose vertices lie in $U$ and whose edges
respect the semantic metric $g$.

Restrictions satisfy sheaf gluing axioms.
\end{definition}

This structure ensures the compatibility of hypergraphs across patches.

\section{Content Atoms}

\subsection{Polyxan atoms revisited}

Polyxan atoms were introduced as sheaf sections in
Chapter~\ref{ch:semantic-manifold}. We now refine the definition.

\begin{definition}[Content Atom]
A content atom is a local hypergraph fragment
\[
a = (V_a, E_a, \tau_a)
\]
supported within a small open neighborhood $U\subseteq\mathcal{M}$.
\end{definition}

The semantic alphabet $\mathcal{A}$ of Chapter~\ref{ch:lsystem} is the set of
geometric equivalence classes of these atoms.

\subsection{Semantic support and stratification}

For an atom $a$ define:
\[
\operatorname{supp}(a) := \pi_H(V_a \cup E_a) \subseteq \mathcal{M}.
\]

Atoms may lie on different semantic strata
$X_\alpha \subseteq\mathcal{M}$, making stratified hypergraph theory essential
for the RSVP--Polyxan coupling.

\section{Hyperedge Algebra and Structural Composition}

\subsection{Tensor products of atoms}

Given two atoms $a$ and $b$ with disjoint semantic supports, define:

\begin{definition}[Atom Tensor Product]
\[
a \otimes b := (V_a\amalg V_b,\; E_a\amalg E_b,\;\tau_a\amalg\tau_b).
\]
\end{definition}

This operation is the combinatorial analogue of parallel semantic composition.

\subsection{Span composition and gluing}

Let $a$ and $b$ be atoms with partially overlapping support.

Define a \emph{hypergraph span}
\[
V_a \leftarrow K \rightarrow V_b
\]
capturing the overlap region $K$.

The composite atom $a\circ b$ is the pushout:
\[
a\circ b := a \amalg_K b.
\]

\begin{proposition}[Existence of Fiberwise Pushouts]
Hypergraph spans admit fiberwise pushouts in $\mathscr{H}(U)$.
\end{proposition}

\begin{proof}[Sketch]
Follows from the existence of pushouts in the category of sets, together with
compatibility conditions for type assignments and semantic embeddings.
\end{proof}

\section{Quines and Self-Referential Structure}

Polyxan hypergraphs support self-replicating atoms via quine edges.

\subsection{Definition of quines}

A quine is a hyperedge $e$ such that:
\[
\mathrm{sig}(\tau(e)) = \tau(e,e_1,\dots,e_k)
\]
for some ordered list of incident edges.

This allows structural recursion.

\subsection{Quine dynamics}

Rewrite rules may produce quines or annihilate them. Quines correspond to:

\begin{itemize}
  \item recursive semantic structures,
  \item persistent content motifs,
  \item fixed points of cognitive rewrite flow.
\end{itemize}

Chapter~\ref{ch:lsystem} established that such fixed points correspond to
semantic attractors.

\section{Semantic Geometry of Hypergraphs}

\subsection{Metric distortion of hyperedges}

Each hyperedge $e=(v_1,\dots,v_k)$ inherits a geometric length
\[
\ell(e) = \sqrt{g(\sigma_V(v_i),\sigma_V(v_j))}
\]
summed across all incident pairs.

High curvature regions of $\mathcal{M}$ induce distortion in hypergraph
structure.

\subsection{Semantic curvature of hypergraphs}

Define:
\[
\mathcal{K}_H :=
\sum_{e\in E}
\bigl(
\mathcal{K}(\sigma_V(v_1))
+ \cdots +
\mathcal{K}(\sigma_V(v_k))
\bigr),
\]
where $\mathcal{K}$ is the semantic curvature from
Chapter~\ref{ch:semantic-manifold}.

This quantity couples directly to RSVP fields.

\section{Typed Families and Hypergraph Indexing}

\subsection{Type universe}

The type system $\mathcal{T}$ is structured as a multi-sorted algebraic theory.

Types may include:
\begin{itemize}
  \item modalities (visual, textual, auditory),
  \item semantic roles (agent, patient, predicate),
  \item generativity types (template, quine, ephemeral),
  \item and RSVP-induced types (inertial, anisotropic, entropic).
\end{itemize}

\subsection{Indexed families}

Each type $t\in\mathcal{T}$ determines a fiberwise family
\[
H_t := \{ x\in H : \tau(x) = t \}.
\]

These families will play a central role in the type-theoretic treatment of
rewrite semantics (Chapter~\ref{ch:ethics}).

\section{Hypergraph Field Coupling}

Typed hypergraphs are not passive structures; they couple to RSVP fields.

\subsection{Coupling to the scalar field}

Define:
\[
\Phi_H := \sum_{v\in V} \Phi(\sigma_V(v)).
\]

This is the accumulated semantic potential of the hypergraph.

\subsection{Coupling to the vector field}

Define a directional flow functional:
\[
\mathbf{v}_H := \sum_{v\in V} \mathbf{v}(\sigma_V(v)).
\]

This encodes anisotropic semantic migration.

\subsection{Coupling to entropy}

Define:
\[
S_H := \sum_{v\in V} S(\sigma_V(v)).
\]

This is the hypergraph’s distributed entropy.

\begin{proposition}[Energy Bound]
The RSVP energy bound (Chapter~\ref{ch:rsvp-fields}) extends to hypergraphs:
\[
\|\nabla \Phi_H\|^2
\le
C_1 + C_2\,|\mathcal{K}_H|.
\]
\end{proposition}

\section{Summary}

This chapter introduced typed Polyxan hypergraphs as the structural layer
beneath semantic content:
\begin{itemize}
  \item hypergraphs fibered over $\mathcal{M}$,
  \item a sheaf of local hypergraphs $\mathscr{H}$,
  \item a type universe $\mathcal{T}$ and associated indexing families,
  \item tensor and span composition of atoms,
  \item quine and recursive structures,
  \item and RSVP field couplings.
\end{itemize}

These objects form the foundation for structural dynamics and semantic
integration. The next chapter develops the categorical and sheaf-theoretic
mechanisms required for higher-level semantic embeddings.
